\section{宣统时期的账本叙事：制度崩解与地方时间}
1909年至1912年的政治转换不能被单一革命中心、退位诏书或后出的胜利叙事垄断。本节分别追踪咨议、铁路、灾荒、军队、各省独立、谈判与退位的材料线索，并保留电报、报刊和地方行动之间可能不同步的时间。
\subsection{1900—1909：事件、文书与地方后果}
这一组节点按账本的日期顺序排列。相邻事件并不必然构成单一因果，却共同显示信息、资源与地方行动如何在同一时段交错。
\hypertarget{ledger:EQ-1909-PROVINCIAL-ASSEMBLIES}{}\paragraph{宣统元年（1909年）}各省咨议局开办，地方精英获得新的公开政治场域。核心取证为《宣统政纪》宣统元年相关条；宫中朱批奏折、内阁大库题本与《清史稿》相关志传。这类上谕、题本或部议首先证明机构处理了何种命令、呈报与责任，而不自动证明地方收到、执行或接受。账本所列条件为：预备立宪需要有限度吸纳地方意见。其后果被限定为：选举资格、代表性和实际权力均受限制。取证保留：以宣统元年（1909年）的同期文书为定位起点；官修记录、奏报或外交文本服务特定机构立场，具体日期、规模、地方执行与对手视角须以独立材料复核。
\hypertarget{ledger:EQ-1909-RAILWAY-RECOVERY-MOVEMENT}{}\paragraph{宣统元年（1909年）}铁路国有化和路权问题引发各省商绅、公司和官府的争论。核心取证为《宣统政纪》宣统元年相关条；户部奏销、盐政、河工或海关档案。这类上谕、题本或部议首先证明机构处理了何种命令、呈报与责任，而不自动证明地方收到、执行或接受。账本所列条件为：外债、地方集资与中央财政需求相冲突。其后果被限定为：“收回路权”诉求包含不同地区和利益群体。取证保留：以宣统元年（1909年）的同期文书为定位起点；官修记录、奏报或外交文本服务特定机构立场，具体日期、规模、地方执行与对手视角须以独立材料复核。
\subsection{1910—1919：事件、文书与地方后果}
这一组节点按账本的日期顺序排列。相邻事件并不必然构成单一因果，却共同显示信息、资源与地方行动如何在同一时段交错。
\hypertarget{ledger:EQ-1910-ZIZHENG-YUAN}{}\paragraph{宣统二年（1910年）}资政院开院，中央层面的立宪咨询机构开始运作。核心取证为《宣统政纪》宣统二年相关条；宫中朱批奏折、内阁大库题本与《清史稿》相关志传。这类上谕、题本或部议首先证明机构处理了何种命令、呈报与责任，而不自动证明地方收到、执行或接受。账本所列条件为：朝廷试图展示宪政进程并控制代表参与。其后果被限定为：咨询权与立法权的界限成为持续争议。取证保留：以宣统二年（1910年）的同期文书为定位起点；官修记录、奏报或外交文本服务特定机构立场，具体日期、规模、地方执行与对手视角须以独立材料复核。
\hypertarget{ledger:EQ-1910-FAMINE-AND-UNREST}{}\paragraph{宣统二年（1910年）}灾荒、物价和财政压力持续影响晚清地方社会与政治动员。核心取证为《宣统政纪》宣统二年相关条；地方志、督抚奏折及地方档案。编年记录适于钉住事件在朝廷文书中的时间位置，却需以地方、对方或物质材料检验其范围和后果。账本所列条件为：自然环境、市场波动和行政能力交互作用。其后果被限定为：地方灾情与政治行动不可简单建立单因果关系。取证保留：以宣统二年（1910年）的同期文书为定位起点；官修记录、奏报或外交文本服务特定机构立场，具体日期、规模、地方执行与对手视角须以独立材料复核。
\hypertarget{ledger:EQ-1911-RAILWAY-NATIONALIZATION}{}\paragraph{宣统三年（1911年）}清廷宣布铁路干线国有化并以外债筹资，保路运动迅速扩大。核心取证为《宣统政纪》宣统三年相关条；宫中朱批奏折、内阁大库题本与《清史稿》相关志传。这类上谕、题本或部议首先证明机构处理了何种命令、呈报与责任，而不自动证明地方收到、执行或接受。账本所列条件为：中央财政、外债与地方商办铁路利益冲突。其后果被限定为：政策公告、地方抗议和后续镇压应分开记录。取证保留：以宣统三年（1911年）的同期文书为定位起点；官修记录、奏报或外交文本服务特定机构立场，具体日期、规模、地方执行与对手视角须以独立材料复核。
\hypertarget{ledger:EQ-1911-SICHUAN-RAILWAY-PROTEST}{}\paragraph{宣统三年（1911年）}四川保路运动升级，赵尔丰处置引发更大政治危机。核心取证为《宣统政纪》宣统三年相关条；地方志、督抚奏折及地方档案。编年记录适于钉住事件在朝廷文书中的时间位置，却需以地方、对方或物质材料检验其范围和后果。账本所列条件为：地方股东、商绅、学生和官府对路权与财政有不同诉求。其后果被限定为：参与规模和暴力过程须依当地档案与报刊核验。取证保留：以宣统三年（1911年）的同期文书为定位起点；官修记录、奏报或外交文本服务特定机构立场，具体日期、规模、地方执行与对手视角须以独立材料复核。
\hypertarget{ledger:EQ-1911-WUCHANG-UPRISING}{}\paragraph{宣统三年（1911年）}武昌起义爆发，新军与革命团体控制武汉部分地区。核心取证为《宣统政纪》宣统三年相关条；军机处档、战事方略及地方奏报。编年记录适于钉住事件在朝廷文书中的时间位置，却需以地方、对方或物质材料检验其范围和后果。账本所列条件为：新军组织、革命网络和铁路危机相互叠加。其后果被限定为：起义的即时参与者、城市控制与后续省份响应需分层叙述。取证保留：以宣统三年（1911年）的同期文书为定位起点；官修记录、奏报或外交文本服务特定机构立场，具体日期、规模、地方执行与对手视角须以独立材料复核。
\hypertarget{ledger:EQ-1911-PROVINCIAL-INDEPENDENCE}{}\paragraph{宣统三年（1911年）}多省宣布脱离清廷，帝国行政和财政网络快速分裂。核心取证为《宣统政纪》宣统三年相关条；宫中朱批奏折、内阁大库题本与《清史稿》相关志传。这类上谕、题本或部议首先证明机构处理了何种命令、呈报与责任，而不自动证明地方收到、执行或接受。账本所列条件为：武昌起义、地方军政力量和立宪派选择共同推动变化。其后果被限定为：“独立”声明不等于各省形成稳定统一控制。取证保留：以宣统三年（1911年）的同期文书为定位起点；官修记录、奏报或外交文本服务特定机构立场，具体日期、规模、地方执行与对手视角须以独立材料复核。
\hypertarget{ledger:EQ-1911-YUAN-SHIKAI-RETURN}{}\paragraph{宣统三年（1911年）}清廷起用袁世凯处理北方军政和议和事务。核心取证为《宣统政纪》宣统三年相关条；宫中朱批奏折、内阁大库题本与《清史稿》相关志传。这类上谕、题本或部议首先证明机构处理了何种命令、呈报与责任，而不自动证明地方收到、执行或接受。账本所列条件为：中央缺乏可依赖的军队与谈判人选。其后果被限定为：袁在清廷、北洋与革命派之间的策略塑造退位进程。取证保留：以宣统三年（1911年）的同期文书为定位起点；官修记录、奏报或外交文本服务特定机构立场，具体日期、规模、地方执行与对手视角须以独立材料复核。
\hypertarget{ledger:EQ-1912-REPUBLIC-FOUNDING}{}\paragraph{宣统四年（1912年）}中华民国临时政府在南京成立，与清廷退位谈判并行。核心取证为《宣统政纪》宣统四年相关条；宫中朱批奏折、内阁大库题本与《清史稿》相关志传。这类上谕、题本或部议首先证明机构处理了何种命令、呈报与责任，而不自动证明地方收到、执行或接受。账本所列条件为：各省独立与南北议和创造新的政权竞争格局。其后果被限定为：临时政府的法理主张与实际控制范围需要区分。取证保留：以宣统四年（1912年）的同期文书为定位起点；官修记录、奏报或外交文本服务特定机构立场，具体日期、规模、地方执行与对手视角须以独立材料复核。
\hypertarget{ledger:EQ-1912-QING-ABDICATION}{}\paragraph{宣统四年（1912年）}宣统帝颁布退位诏书，清帝国终结。核心取证为《宣统政纪》宣统四年相关条；宫中朱批奏折、内阁大库题本与《清史稿》相关志传。这类上谕、题本或部议首先证明机构处理了何种命令、呈报与责任，而不自动证明地方收到、执行或接受。账本所列条件为：南北谈判、财政军政压力与皇室待遇安排共同促成退位。其后果被限定为：退位文本规定的优待与地方政治现实须分别追踪。取证保留：以宣统四年（1912年）的同期文书为定位起点；官修记录、奏报或外交文本服务特定机构立场，具体日期、规模、地方执行与对手视角须以独立材料复核。
